\section{Introduction}
\label{sec:intro}

Generative super-resolution (SR) increasingly runs in a
\emph{latent} space: an autoencoder compresses images, a diffusion
model denoises the high-resolution (HR) latent, and a decoder returns
pixels~\citep{rombach2022high,saharia2023image,wang2024exploiting}.
The low-resolution (LR) observation enters as a condition, typically
the encoded upsampled LR image $\zlr=\Enc(\xup)$.
This design assumes a convenient causal story:
if $\zlr$ is a better representation of the HR target, the reverse
process should emit a better $\zsr$, and SR metrics should follow.

That story is not a theorem.
It confuses two jobs that share a decoder but not an objective
(\Cref{tab:two-jobs}):
soft-decode $\Dec(\zlr)$ scores what is already in the LR code;
\lsrmethod{} $\Dec(\zsr)$ scores what a conditional DDPM uses while
sampling $p(\zhr\mid\zlr)$.

This paper is a controlled diagnosis of that transfer, not a new
SOTA SR method.
We freeze the data, the UNet skeleton, the noise schedule, and the
evaluation noise, and change one link at a time
(\Cref{fig:causal-chain}): representation (\vaesr), injection
(concat vs.\ spatial FiLM), reverse-chain measurement, then
inference-time LR-consistency guidance.
Headline outcomes---a large soft-decode gain that concat and FiLM
do not consume, a mid-reverse copy that later collapses, and late
guidance that recovers about half the gap---are in~\S\ref{sec:experiments}.
The intervention that works is not the one the representation
hypothesis suggested.

\paragraph{Scope.}
No new backbone, tokenizer, or trained guidance network.
The FiLM null is for this conditioner on this frozen \vaesr{} code.
$\Dec(\zlr)$ is not an SR method.
Citation hygiene:~\S\ref{sec:setup-hygiene}.

\begin{figure*}[t]
  \centering
  \resizebox{\textwidth}{!}{%
  \begin{tikzpicture}[
      font=\small,
      box/.style={draw, rounded corners=2pt, align=center, inner sep=6pt, minimum height=1.35cm},
      arr/.style={-{Latex}, thick}
    ]
    \node[box, fill=blue!8] (rep) {%
      \textbf{1. Representation}\\[2pt]
      \vaeone{} $\to$ \vaesr\\
      $\Dec(\zlr)$: $26.15\to 28.48$\,dB\\
      \textit{code gets better}};
    \node[box, fill=orange!12, right=0.55cm of rep] (inj) {%
      \textbf{2. Injection}\\[2pt]
      concat or spatial FiLM\\
      SR: $26.26\to 26.48$\,dB\\
      FiLM $\Delta{=}-0.03$ (null)\\
      \textit{reader is not the hole}};
    \node[box, fill=yellow!20, right=0.55cm of inj] (traj) {%
      \textbf{3. Reverse chain}\\[2pt]
      $\cos(\zhat,\zlr)$\\
      copy $0.985$ at $t{=}656$\\
      $t{=}0$: $0.786$\\
      \textit{late HR prior}};
    \node[box, fill=green!12, right=0.55cm of traj] (guid) {%
      \textbf{4. Guidance}\\[2pt]
      DPS on $t\le 500$\\
      $\lamg{=}200$: $26.33\to 27.50$\\
      $56\%$ of $\approx 2$\,dB gap\\
      \textit{no new training}};
    \draw[arr] (rep) -- (inj);
    \draw[arr] (inj) -- (traj);
    \draw[arr] (traj) -- (guid);
  \end{tikzpicture}%
  }
  \caption{Causal chain.
  Each box changes one link and holds the others fixed.
  Numbers are the locked outcomes of~\S\ref{sec:experiments}
  ($n{=}64$ except the last box, which is the $n{=}256$ dose).}
  \label{fig:causal-chain}
\end{figure*}
